% FILE: 00_project_identity.tex
% Project: comparisons
% Role: cross-system-capability-comparison-corpus
% Protocol: v30.1.1

\section{Project Identity}
\label{sec:comparisons-identity}

\subsection{Purpose}
\label{sec:comparisons-identity-purpose}

The \texttt{comparisons} corpus is the cross-system capability authority for the
process-intelligence research program. Its purpose is to establish, with
source-traceable evidence, the precise capability boundaries of three distinct
execution stacks: PM4Py (CPython/Pandas), wasm4pm (WebAssembly/Rust execution
engine), and wasm4pm-compat (Rust compile-time type-law enforcement layer). The
corpus does not argue that one stack is universally superior. It establishes what
each stack can and cannot do, at what cost, and at which enforcement level.

The central architectural thesis is that PM4Py is a data-science prototyping
tool with runtime-only validation; wasm4pm is an edge-native execution engine with
strong Rust types but no admission gates; and wasm4pm-compat is the type-law
authority that makes invalid conformance metrics (fitness $> 1.0$), unvalidated
event log input, and generic exception handling compile-time impossibilities rather
than runtime defects. This three-way comparison is the evidentiary basis for
design decisions across all downstream engineering surfaces in the research
program.

\subsection{Repository Surface}
\label{sec:comparisons-identity-surface}

The corpus resides at \texttt{/Users/sac/process-intelligence/comparisons/} and
consists of eleven substantive markdown documents authored under the v30.1.1
research protocol, all timestamped 2026-05-31 to 2026-06-01. The eleven documents
cover: architecture comparison, conformance metric surfaces, algorithm coverage
matrix, type safety comparison, admission/refusal law comparison, WASM vs.\ native
performance benchmarks, declarative vs.\ imperative conformance, lattice vs.\ trace
compliance checking, PM4Py vs.\ wasm4pm conformance details, algorithm comparison,
and core lifecycle algorithms.

The ALIVE gate assessment (\texttt{checkpoints/ALIVE\_GATE\_ASSESSMENT.md})
explicitly lists the comparisons file count as criterion 8: target $\geq 4$,
actual 11, verdict MET. The corpus thus satisfies one of the twelve criteria that
sealed the PROCESS\_INTELLIGENCE\_ALIVE\_001 checkpoint on 2026-05-31.

\subsection{Primary Research Role}
\label{sec:comparisons-identity-role}

Within the dissertation, the comparisons corpus performs the role of
\emph{cross-system capability comparison corpus}. It is the empirical surface on
which claims about type-law superiority, performance parity, and algorithmic gap
structure are grounded. No downstream engineering claim about wasm4pm-compat
design choices can be made without a comparisons-corpus entry tracing it.

The corpus also provides the comparative baseline for the M\&A claim layer: the
\texttt{PM4PY\_ORACLE\_RECEIPT} and \texttt{WASM4PM\_GAP\_RECEIPT} in the receipt
registry are explicitly described as covering the comparison surface. Any
board-level fitness, precision, or algorithmic coverage claim must trace to a
comparisons entry.

\subsection{Key Artifacts}
\label{sec:comparisons-identity-artifacts}

The key artifacts produced by or consumed in this corpus are:

\begin{itemize}
  \item \texttt{pm4py\_vs\_wasm4pm\_architecture.md} --- Three-system architectural
    comparison with mathematical WF-net soundness definitions and LTL semantics.
  \item \texttt{CONFORMANCE\_METRIC\_SURFACES.md} --- Source-traced demonstration
    that \texttt{Metric<KIND, NUM, DEN>} with \texttt{Require<\{NUM <= DEN\}>:
    IsTrue} makes out-of-range quality metrics a compile-time error.
  \item \texttt{ALGORITHM\_COVERAGE\_MATRIX.md} --- Eleven-row coverage table
    classifying every algorithm as Full, Shape, Partial, or No across the three
    stacks.
  \item \texttt{TYPE\_SAFETY\_COMPARISON.md} --- The laundering test: can
    unvalidated event logs reach algorithm results without detection? The answer
    is No only in wasm4pm-compat.
  \item \texttt{ADMISSION\_REFUSAL\_VS\_PM4PY\_ERRORS.md} --- Named-law refusal
    architecture vs.\ Python exception strings.
  \item \texttt{wasm\_vs\_native\_pm\_performance.md} --- Performance figures:
    8.2M events/sec wasm4pm vs.\ 6.3K events/sec PM4Py for alignment conformance;
    $< 1$ms WASM instantiation vs.\ 10--50ms native.
  \item \texttt{lattice\_vs\_trace\_compliance\_checking.md} --- Formal proof that
    the witness join-semilattice $(W, \sqsubseteq, \sqcup, \bot, \top)$ with
    Monotonicity Axiom 2 is sound and complete with respect to trace-based
    language acceptance.
\end{itemize}
